\TLEFigure{2}{À gauche : spectre d'émission d'un sylphe. À droite : coefficient de transmission de l'atmosphère avec deux taux d'humidité $20\%$ et $80\%$. Ces mesures ont été réalisées par un détecteur au sol placé sous une source située à $80\text{ km}$ d'altitude -- Figures tirées de l'article \textit{Spectrum of Red sprites}, G. Milikh, J. A. Valdiviab \& K. Papadopoulos, Journal of Atmospheric and Solar-Terrestrial Physics, Vol. 60, p. 907-915, 1998}

Du fait de la rotondité de la Terre de rayon $R_T$, on s'interroge sur la possibilité de vision de sylphes au-dessus des Alpes.

\TLEFigure{3}{Photos tirées de l'article \textit{Sylphes rouges, Jets Bleus et elfes}, E. Blanc, R. Roche, T. Farges \& F. Simonet publié dans la revue \textsc{Chocs}, Revue scientifique et technique de la Direction des applications militaires du CEA, numéro 26, décembre 2002}

\begin{itemize}
    \item[--] \textbf{2.} La distance $d$ qui sépare à vol d'oiseau le Pic du Midi de Bigorre de la chaîne des Alpes françaises est-elle de l'ordre de 500, 1000 ou 1500 km ?
    
    Sans prendre en compte l'altitude du Pic du Midi de Bigorre, tracer sur un schéma approprié la ligne d'horizon qui en part et qui passe au dessus des Alpes à une altitude $h$. En faisant les hypothèses qui s'imposent exprimer $h$ en fonction de $R_T$ et $d$ et calculer sa valeur numérique.
    
    Est-il possible de voir des sylphes alpins sans effet de réfraction atmosphérique depuis le Pic du Midi de Bigorre ? Peut-on voir le Mont-Blanc par beau temps ?
\end{itemize}

Afin de pouvoir observer régulièrement des sylphes et de développer des recherches visant à mieux comprendre leur origine, on utilise depuis 2001 sur la station spatiale internationale (\textsc{iss}), deux microcaméras fixées sur un hublot. Les deux caméras observent selon la verticale. L'une de ces deux caméras opère dans le visible tandis que l'autre est équipée d'un filtre de très faible bande passante centré à $761\text{ nm}$. Sur la figure 4, on peut observer les deux images obtenues lors d'un enregistrement. Cet enregistrement a été réalisé lors d'un orage important.
